\section{Recent Applications in Document OCR}

The practical relevance of optical character recognition extends well beyond academic benchmarks, with deployed OCR systems now underpinning a wide range of real-world document processing workflows across industries and public institutions. The maturation of deep learning-based detection and recognition architectures has enabled OCR to move from controlled laboratory settings into operational environments characterized by document diversity, scale requirements, and strict accuracy demands. This section surveys three representative application domains that illustrate the breadth of contemporary OCR deployment: administrative and financial document processing, identity document recognition, and historical archive transcription. Each domain presents a distinct set of technical challenges that contextualize the practical value of the Vietnamese OCR system developed in this thesis.

\subsection{Administrative and Financial Document Processing}
The digitization of administrative and financial documents, including invoices, receipts, purchase orders, and tax forms, represents one of the most commercially mature applications of OCR technology. Organizations processing large volumes of such documents have strong economic incentives to automate extraction of structured information such as vendor names, dates, line items, and totals, reducing reliance on manual data entry and accelerating downstream business processes. Early OCR-based approaches to this problem relied on template matching against fixed document layouts, but the diversity of invoice formats across vendors and jurisdictions made template-based systems brittle and expensive to maintain \cite{Palm2017}.

More recent work has framed invoice and receipt understanding as a joint OCR and information extraction problem, in which text detection and recognition are coupled with named entity recognition and document layout analysis to produce structured output from unstructured document images. The CORD dataset \cite{Park2019} and the SROIE benchmark \cite{Huang2019} established standardized evaluation frameworks for receipt OCR and key information extraction, driving the development of end-to-end trainable systems that simultaneously localize text, transcribe it, and classify it into semantic categories. LayoutLM \cite{Xu2020} extended transformer-based language models to the document understanding domain by incorporating two-dimensional positional embeddings that encode the spatial coordinates of each text token alongside its semantic embedding, enabling the model to exploit document layout structure as an additional signal for information extraction. Subsequent iterations including LayoutLMv2 \cite{Xu2022} and LayoutLMv3 \cite{Huang2022} further integrated visual features from document images directly into the pretraining objective, producing unified multimodal representations that have become the dominant paradigm for structured document understanding

The relevance of this application domain to Vietnamese OCR is direct. Vietnamese government agencies and enterprises generate large volumes of administrative documents including tax invoices, procurement records, and official correspondence, all of which are candidates for automated digitization. The orthographic complexity of Vietnamese text imposes additional demands on the OCR component of such pipelines, as recognition errors on diacritical marks can produce semantically distinct output that corrupts downstream information extraction.

\subsection{Identity Document Recognition}
Identity document recognition, encompassing the automated extraction of personal information from national identity cards, passports, driving licenses, and similar credentials, constitutes another high-value OCR application domain with stringent accuracy and reliability requirements. Errors in recognized identity fields can have serious consequences in contexts such as border control, financial onboarding, and voter registration, making this one of the most demanding real-world OCR deployment scenarios.

Early identity document OCR systems exploited the highly structured layouts of standardized documents, using template registration to locate predefined field regions before applying OCR within each region. The introduction of machine-readable zones (MRZ) on passports and identity cards provided a constrained recognition target amenable to classical OCR techniques, given the limited character set and fixed typography of MRZ fields. However, the diversity of identity document formats across countries and the presence of security printing features, holographic overlays, and variable background textures have motivated the adoption of deep learning-based approaches that generalize across document types without requiring per-format template engineering. The MIDV-500 dataset \cite{Bulatov2019} established a standardized benchmark for identity document analysis by providing annotated samples across 50 document types from different countries under varied capture conditions including flat, curved, and perspective-distorted document images, enabling more rigorous evaluation of generalizable identity document recognition systems.

More recent work has framed identity document recognition as an end-to-end structured information extraction problem, combining text detection, recognition, and field classification within a unified pipeline \cite{Tran2019}. Rather than treating each processing stage independently, these systems jointly optimize detection and recognition objectives, enabling the model to leverage layout structure as a supervisory signal for improving transcription accuracy across variable document formats. For Vietnamese identity documents specifically, the transition from the old nine-digit citizen identification card to the new twelve-digit CCCD format introduced new layout conventions and font specifications that required OCR systems to be retrained or fine-tuned, highlighting the importance of adaptable recognition architectures \cite{Nguyen2022}. The cascaded detection-recognition design adopted in this thesis is well suited to this setting, as the modular architecture allows the detection and recognition components to be independently fine-tuned when document specifications change.

Across both application domains surveyed in this section, a consistent pattern emerges: the performance of deployed OCR systems is ultimately bounded by the quality of the underlying text detection and recognition components, and improvements to these core components translate directly into downstream application quality. This observation reinforces the motivation for the thesis at hand, which targets the detection and recognition pipeline as the primary locus of improvement for Vietnamese document OCR.
